% Generated by scripts/txt2tex.py from chapters/chapter_01.txt. Edit the TXT, not this file.
\chapterquestion{מי כתב את התשובה שמופיעה עכשיו על המסך, אם שום אדם לא ראה את השאלה?}
\chapter{מהי בינה מלאכותית?}\label{ch:01}
\section{רגע לפני Enter}
דמיינו שעת ערב מאוחרת. הבצק שהקוראת השאירה על השיש לא תפח, והיא פותחת חלון של מערכת בינה מלאכותית ומקלידה: "למה הבצק שלי לא תופח, ומה אפשר לעשות עכשיו?" כעבור שניות מתחילות להופיע על המסך מילים, אחת אחרי השנייה. לתמונה הזאת נחזור לאורך הספר.
אף מתכנת לא כתב מראש את התשובה הזאת. איש לא ניחש שהקוראת תשאל דווקא על בצק, דווקא הערב, דווקא בניסוח הזה. ובכל זאת התשובה מנוסחת, נוגעת בעניין, ולפעמים אפילו נכונה. מאיפה היא הגיעה?
הספר יוצא למסע אל מאחורי הקלעים של הרגע הקטן הזה.
\section{שם שנולד בהצעת מחקר}
ב-1950 פתח טיורינג מאמר בשאלה "Can machines think?", כלומר: האם מכונות יכולות לחשוב? הוא הציע ניסוי שכינה "משחק החיקוי": חוקר בחדר נפרד שואל שאלות ומנסה לקבוע מי מבין שני משתתפים סמויים הוא הגבר ומי האישה. ואז טיורינג שואל: מה יקרה אם מכונה תתפוס את מקומו של אחד מהם \cite{S-ch01-02}.
השם עצמו נולד כעבור חמש שנים, בהצעת מחקר מ-31 באוגוסט 1955. ארבעה חוקרים, מקארתי, מינסקי, רוצ'סטר ושאנון, הציעו לקיים בקיץ 1956 מחקר של חודשיים בהשתתפות עשרה אנשים, ונתנו לנושא את השם artificial intelligence, כלומר בינה מלאכותית \cite{S-ch01-01}. ההצעה נשענה על השערה נועזת: שכל היבט של למידה, וכל תכונה אחרת של אינטליגנציה, אפשר עקרונית לתאר בדיוק רב כל כך, שמכונה תוכל לדמות אותו. בין שבעת ההיבטים של הבעיה שההצעה מנתה הופיע כבר אז "שיפור עצמי" (Self-Improvement) \cite{S-ch01-01}.
\section{הגדרה שמסרבת להתקבע}
עשרות שנים אחר כך עדיין אין הגדרה מדויקת ומוסכמת, כפי שמודה פאנל המומחים של פרויקט AI100 בסטנפורד (2016) \cite{S-ch01-07}. הפאנל מביא את ההגדרה של נילס נילסון (Nils J. Nilsson): בינה מלאכותית היא פעילות שמטרתה להפוך מכונות לאינטליגנטיות, ואינטליגנציה היא התכונה שמאפשרת לישות לתפקד כראוי, ומתוך ראייה קדימה, בסביבה שלה \cite{S-ch01-07}. הפאנל מתאר גם את "אפקט ה-AI" (AI effect): כשאנשים מתרגלים לטכנולוגיה, היא מפסיקה להיחשב בינה מלאכותית \cite{S-ch01-07}.
הארגון לשיתוף פעולה ולפיתוח כלכלי (OECD) בחר בניסוח מעשי יותר. לפי ההגדרה הרשמית שלו, מערכת בינה מלאכותית היא מערכת מבוססת מכונה שמסיקה מהקלט שהיא מקבלת, לשם מטרות מפורשות או משתמעות, כיצד לייצר פלטים כמו תחזיות, תוכן, המלצות או החלטות, שעשויים להשפיע על סביבות פיזיות או וירטואליות \cite{S-ch01-09}. בעדכון ההגדרה (2023–2024) נוספה המילה "תוכן", כדי להבהיר שגם מערכות של בינה מלאכותית יוצרת (generative AI), כלומר מערכות שמפיקות טקסט ותוצרים חדשים, נכללות בה \cite{S-ch01-10}.
ובערב של הקוראת: הקלט הוא שאלה על בצק, הפלט הוא תוכן על המסך, ובאמצע עומדת המילה החשובה, "מסיקה". ההגדרה לא שואלת אם המכונה חושבת, אלא מה היא עושה עם מה שהיא מקבלת.
\section{מתכון מול טעימה}
תוכנה רגילה דומה למתכון כתוב: מתכנת כותב לכל מצב הוראה מפורשת, אם קרה כך, עשה כך. זה עובד כל עוד אפשר לנסח את הכללים. נסו לכתוב רשימת כללים שמזהה את הספרה 7 בכתב יד: עם קו באמצע או בלי, זקופה או נטויה, מעוגלת או חדה. כל כלל חדש מוליד חריג חדש.
טום מיטשל ניסח את ההבדל בחיבור משנת 2006: מדעי המחשב עסקו בעיקר בשאלה איך לתכנת מחשבים ביד, ואילו למידת מכונה (machine learning) שואלת איך לגרום למחשבים לתכנת את עצמם, מתוך ניסיון ומעט מבנה התחלתי \cite{S-ch01-03}. לדבריו, גם בזיהוי דיבור וגם בראייה ממוחשבת, מערכות מאומנות מדויקות יותר מתוכנות שנכתבו ביד. ונכון ל-2006, יותר מ-85\% מהדואר בכתב יד בארצות הברית מוין אוטומטית, בתוכנה שאומנה בלמידת מכונה \cite{S-ch01-03}.
עצרו רגע מול המספר הזה. על כל מעטפה כזאת כתב מישהו בכתב ידו, עם הנטייה והעיגולים שרק לו יש. איש לא מנה בפני המכונה את כל הדרכים לכתוב כל אות וכל ספרה, ובכל זאת היא למדה לפענח אותן מתוך דוגמאות. מיומנות שקשה להסביר במילים נלמדה בלי הסבר.
אם תוכנה רגילה היא מתכון, מערכת לומדת דומה לטבח שמעולם לא קרא מתכון, אבל טעם אינספור תבשילים ושמע בכל פעם מה הצליח ומה לא. הכללים שלו לא נכתבו בשום מקום. הם התגבשו.
טיורינג חשב בכיוון דומה. הוא הציע לבנות מכונה שמדמה מוח של ילד ולא של מבוגר, ותיאר את מוח הילד כמחברת חדשה מחנות כלי כתיבה: מעט מנגנון והרבה דפים ריקים \cite{S-ch01-02}.
\section{מה הופך מערכת ל"לומדת"}
\begin{conceptbox}{למידה לפי משימה, מדד וניסיון}
מיטשל מציע הגדרה בת שלושה רכיבים: מכונה לומדת ביחס למשימה T (task), מדד ביצוע P (performance) וסוג ניסיון E (experience), אם היא משפרת באופן אמין את הביצועים שלה במדד P במשימה T בעקבות הניסיון E \cite{S-ch01-03}.
לדוגמה, מסנן דואר זבל:
המשימה (T): להבחין בין הודעה רצויה להודעת זבל.
המדד (P): החלק של ההודעות שסווגו נכון.
הניסיון (E): הודעות קודמות שמשתמשים כבר סימנו כ"זבל" או כ"תקין".
אם המסנן טועה פחות ככל שהוא רואה יותר הודעות מסומנות, הוא לומד.
\end{conceptbox}
ההגדרה לא דורשת הבנה או תודעה, רק שיפור שאפשר למדוד. מיטשל אף רואה בלמידת מכונה צאצא טבעי של המפגש בין מדעי המחשב לסטטיסטיקה \cite{S-ch01-03}: כדי לדעת אם משהו השתפר, צריך לספור ולהשוות.
\section{שלוש בובות, אחת בתוך השנייה}
כדי לסדר את שלושת המושגים, הספר הזה מציע תמונה משלו: בובת מטריושקה. הבובה החיצונית היא בינה מלאכותית: מערכות שמקבלות קלט ומסיקות ממנו פלט. בתוכה יושבת למידת מכונה: מערכות שההתנהגות שלהן נבנית מתוך ניסיון ולא נכתבת כולה ביד. ובפנים, הקטנה ביותר: למידה עמוקה (deep learning). לא כל בינה מלאכותית לומדת, ולא כל למידה היא עמוקה.
מומחי AI100 מונים כמה גורמים שהזינו את מהפכת הבינה המלאכותית, ובראשם הבשלתה של למידת המכונה \cite{S-ch01-08}. את הלמידה העמוקה הם מתארים כסוג של רשתות נוירונים מלאכותיות (artificial neural networks), שמאומנות בשיטה שנקראת התפשטות לאחור (backpropagation) \cite{S-ch01-08}. רשת נוירונים היא מבנה חישובי של יחידות פשוטות רבות שמחוברות זו לזו.
ומה פירוש "עמוקה"? לפי סקירה שפרסמו לקון, בנג'יו והינטון ב-Nature ב-2015, מערכת כזאת בנויה משכבות עיבוד רבות, וכל שכבה בונה את התיאור שלה לנתונים מתוך התיאור של השכבה שלפניה. כך נוצרות כמה רמות של הפשטה \cite{S-ch01-05}. לתיאור כזה קוראים ייצוג (representation). זה דומה לילד שלומד לקרוא: קודם קווים, אחר כך אותיות, מילים ולבסוף משפט שלם.
ומה קובע את התיאור שכל שכבה מפיקה? אוסף גדול של מספרים פנימיים, הפרמטרים, שאפשר לדמות לכפתורים במערבל קול; בפרק 3 נפתח רשת כזאת ונפגוש אותם בשמות משקולות והטיות. התפשטות לאחור היא האלגוריתם שמסמן למכונה איך לשנות את הכפתורים האלה בכל שכבה \cite{S-ch01-05}. דמיינו טכנאי שמשווה את מה שנשמע למה שהיה אמור להישמע, ומזיז כל כפתור מעט בכיוון שמקטין את הפער. לפי הסקירה, השיטות האלה שיפרו מאוד את התוצאות הטובות ביותר שהושגו עד אז בזיהוי דיבור, בזיהוי עצמים בתמונות ובתחומים רבים אחרים \cite{S-ch01-05}.
\begin{figure}[htbp]
\centering
% Figure 1.1: three nested circles (AI > ML > DL). Book framework, no external data.
% Gold is used only for marker dots; label text stays Midnight Blue (contrast, R-ch01-16).
\begin{bookdiagram}
  \fill[midnight!7,draw=midnight,line width=1.2pt] (0,0) circle (3.3);
  \fill[turquoise!14,draw=turquoise,line width=1.2pt] (0,-0.95) circle (2.3);
  \fill[violet!16,draw=violet,line width=1.2pt] (0,-1.95) circle (1.3);
  \node[font=\headingfont\normalsize] at (0,2.55) {בינה מלאכותית};
  \node[font=\headingfont\normalsize] at (0,0.75) {למידת מכונה};
  \node[font=\headingfont\normalsize,align=center] at (0,-1.95) {למידה\\עמוקה};
  \draw[midnight,line width=0.6pt] (-1.75,2.55) -- (-4.3,2.3);
  \draw[turquoise,line width=0.6pt] (-1.45,0.75) -- (-4.3,0.5);
  \draw[violet,line width=0.6pt] (-0.75,-1.95) -- (-4.3,-1.5);
  \foreach \y in {2.3,0.5,-1.5}{\fill[gold] (-4.3,\y) circle (1.6pt);}
  \node[anchor=east,text width=4.1cm,align=right,font=\footnotesize] at (-4.45,2.3) {מסיקה מתוך הקלט שהיא מקבלת איך לייצר פלט};
  \node[anchor=east,text width=4.1cm,align=right,font=\footnotesize] at (-4.45,0.5) {ההתנהגות שלה נבנית מתוך ניסיון};
  \node[anchor=east,text width=4.1cm,align=right,font=\footnotesize] at (-4.45,-1.5) {ייצוגים שנלמדים בשכבות רבות};
\end{bookdiagram}

\caption{שלוש בובות, אחת בתוך השנייה}\label{fig:1.1}
\figurenote{מקור: מסגרת הספר (לא מקור חיצוני)}
\end{figure}
\begin{thinkbox}
הפאנל של סטנפורד מזכיר שמחשבון אלקטרוני פשוט מחשב מהר בהרבה ממוח אדם, ורואה באינטליגנציה רצף: ההבדל בין מחשבון למוח הוא בקנה מידה, במהירות, באוטונומיה ובכלליות, ולא בסוג \cite{S-ch01-07}.
האם המחשבון בטלפון שלכם הוא בינה מלאכותית? ומה לגבי ההשלמה האוטומטית במקלדת? היכן הייתם מותחים את הקו, והאם הוא יישאר שם בעוד עשור?
\end{thinkbox}
\section{מפת הדרך: מה מחכה מאחורי המסך}
השאלה על הבצק נשלחה. אלה התחנות שבדרך: פרק 2, "מנתונים ללמידה": הדוגמאות שמהן המערכת למדה, אימון, בדיקה והטיות. פרק 3, "הרשת הנוירונית מבפנים": איך מספרים ומשקולות מצטרפים לתחזית. פרק 4, "כיצד מודל שפה מייצר תשובה": איך התשובה נבנית חלק אחר חלק. פרק 5, "בודקים מודלים בעזרת נתונים": איך מבחינים בין הצלחה לאשליה שלה. פרק 6, "מה AI יודע — ומה לא": טעויות בטוחות בעצמן, פרטיות וזכויות. פרק 7, "האדם במעגל והעתיד": תפקידה של הקוראת מול התשובה. נספח, "כך נוצר הספר שאתם קוראים": מאחורי הקלעים של הספר עצמו.
אבל המסע מתחיל הרבה לפני שהקוראת הקלידה: בדוגמאות שמהן המערכת למדה.
\begin{insightsbox}
\begin{itemize}
\item "בינה מלאכותית" היא משפחה של הגדרות; המעשיות שבהן שואלות מה המערכת עושה עם הקלט, ולא אם היא "חושבת" \cite{S-ch01-07,S-ch01-09}.
\item בתוכנה רגילה אדם מנסח את הכללים; במערכת לומדת המחשב "מתכנת את עצמו" מתוך ניסיון ומעט מבנה התחלתי \cite{S-ch01-03}.
\item לפי ההגדרה של מיטשל, מערכת למדה אם הביצועים שלה במשימה, לפי מדד מוגדר, משתפרים באופן אמין בעקבות ניסיון \cite{S-ch01-03}.
\item למידה עמוקה לומדת שכבה על גבי שכבה של ייצוגים \cite{S-ch01-05}. במטריושקה של הספר היא רק הבובה הפנימית, ולא שם נרדף לבינה מלאכותית.
\item הקו בין "כלי" ל"בינה" זז עם ההרגל: מה שמפליא היום ייראה מחר כמו עוד פונקציה בטלפון \cite{S-ch01-07}.
\end{itemize}
\end{insightsbox}
